\documentclass[10pt,letterpaper]{article}

% ---------- Idioma y codificación ----------
\usepackage[spanish]{babel}
\usepackage[T1]{fontenc}
\usepackage[utf8]{inputenc}

% ---------- Márgenes ----------
\usepackage[margin=1.25cm]{geometry}

% ---------- Tipografía ----------
\usepackage{lmodern}

% ---------- Utilidades ----------
\usepackage{xcolor}
\usepackage[hidelinks]{hyperref}
\usepackage{enumitem}
\usepackage{titlesec}
\usepackage{multicol}

% ---------- Espaciado y listas ----------
\setlength{\parindent}{0pt}
\setlength{\parskip}{0.05\baselineskip}
\setlist[itemize]{noitemsep, topsep=0pt, leftmargin=1.2em, after=\vspace{0.05\baselineskip}}

% ---------- Formato de secciones ----------
\titleformat{\section}{
  \large\bfseries\uppercase
}{}{0em}{}[\titlerule]

\titlespacing{\section}{0pt}{3pt}{1pt}

\newcommand{\entry}[4]{
  \begin{minipage}[t]{\linewidth}
    \textbf{#1} \hfill \textit{#2} \\
    \textbf{\textit{#3}} \hfill #4
  \end{minipage}
}

\begin{document}

\pagenumbering{gobble}

\begin{center}
	{\huge\bfseries LUIS ÁNGEL RAMÍREZ VILLANUEVA}\\[1pt]
	Oaxaca de Juárez \textbullet{} abierto a remoto México \textbullet{} (+52) 951 765 2254 \textbullet{}
	\href{mailto:luisrmrz1699@gmail.com}{luisrmrz1699@gmail.com}\\[1pt]
	\href{https://github.com/luisrmrz16}{GitHub} \textbullet{}
	\href{https://www.linkedin.com/in/luisrmrz16/}{LinkedIn}
\end{center}

\section*{Perfil profesional}

Más de 3 años en soporte en sitio y remoto (L1 y algo de L2): mesa de ayuda, tickets, resolución de incidencias y empatía con usuarios. De forma secundaria, mejora herramientas internas operativas con apoyo de IA (``vibecoding'') y nociones básicas de desarrollo. Busca crecer en soporte, service desk o soporte en campo/remoto en México.

\section*{Experiencia laboral}

\entry
{Encargado del Área de Sistemas de Sucursal}
{Oaxaca, México}
{Herst Electronics}
{Junio 2026 --- Actualidad}
\begin{itemize}
	\item Soporte en sitio y mesa de ayuda (Romasa): tickets, Windows 10/11, impresoras, redes y Microsoft Office / Outlook
	\item Apoyo administrativo con Sicofi (facturación), ERP en MS Access y SQL Server básico; bitácora / documentación operativa
	\item Respaldo de BD del ERP y verificación de procesos en servidores; mejora de herramienta interna de inventario (Angular + Supabase) con IA
\end{itemize}

\entry
{Auxiliar de Sistemas}
{Oaxaca, México}
{Herst Electronics}
{Julio 2025 --- Junio 2026}
\begin{itemize}
	\item Soporte en sitio a Romasa y clientes externos: hardware, software, impresoras e incidencias de red
	\item Instalación y preparación de equipos Windows 10/11 (Office y software básico); mantenimiento preventivo de PCs e impresoras
	\item Contribución inicial a sistema interno de inventario y mantenimientos para el área de soporte
\end{itemize}

\entry
{Ingeniero de Soporte en Sitio}
{Oaxaca, México}
{Kenos Scanda}
{Enero 2023 --- Febrero 2025}
\begin{itemize}
	\item Soporte en sitio (L1/L2) para Mobility ADO: oficinas, terminales y sitios operativos; atención a usuarios y tickets
	\item Ciclo de vida de tickets en IBM Maximo y ServiceNow (ITSM): registro, seguimiento, resolución y escalamiento
	\item Active Directory (unión a dominio, validación de usuarios); Office / Outlook en entornos corporativos
	\item Instalación e inventario de equipos; IP fijas, redes / DNS y registro de responsivas
\end{itemize}

\entry
{Técnico de Soporte y Reparación de Equipos de Cómputo}
{Oaxaca de Juárez, Oaxaca}
{Innovaciones Tecnológicas}
{Septiembre 2018 --- Septiembre 2022}
\begin{itemize}
	\item Diagnóstico, reparación y mantenimiento de laptops y PCs; instalación de Windows y macOS; atención a clientes
\end{itemize}

\section*{Proyectos}

\begin{itemize}
	\item \textbf{Herramienta interna de inventario y mantenimientos (Angular + Supabase)} --- Mejorada en Herst como apoyo al soporte operativo; interfaz web con consumo de datos; vibecoding con IA (no portafolio de desarrollador).
\end{itemize}

\section*{Educación}

\entry
{Ingeniería en Sistemas Computacionales (en curso)}
{Oaxaca de Juárez, Oaxaca}
{Instituto de Estudios Universitarios}
{Octubre 2023 --- Enero 2027}

\entry
{Bachillerato General}
{Oaxaca de Juárez, Oaxaca}
{Colegio de Bachilleres del Estado de Oaxaca}
{Enero 2017 --- Diciembre 2019}

\begin{multicols}{2}

	\textbf{\MakeUppercase{Habilidades técnicas}}
	\begin{itemize}
		\item Soporte en sitio/remoto; mesa de ayuda y tickets
		\item Windows 10/11, macOS, Linux; Active Directory
		\item Hardware, impresoras, redes (IP/DNS)
		\item Microsoft Office, Outlook; colaboración
		\item ServiceNow, IBM Maximo (ITSM); SQL Server básico
		\item Bitácora y documentación operativa
		\item Vibecoding con IA; Angular/Supabase básico; Git básico
	\end{itemize}

	\columnbreak

	\textbf{\MakeUppercase{Certificaciones e idiomas}}
	\begin{itemize}
		\item CCNA Routing and Switching (Cisco NetAcad)
		\item IT Essentials (Cisco NetAcad)
		\item Curso Profesional de Git y GitHub (Platzi)
		\item Fundamentos de Ingeniería de Software (Platzi)
		\item Introduction to Generative AI (Google Cloud)
		\item Inglés – lectura técnica; conversación básica
	\end{itemize}

\end{multicols}

\end{document}
